% Abhakseite „Das kann ich" – Handform, nach Einheiten und Verfahrensüberschriften
\newcommand{\vgruppe}[1]{\abhakgruppe{\normalfont\itshape #1}}
\begin{abhakseite}
\abhakgruppe{Einheit 1 · Proportionale Funktion}
\vgruppe{Wertetabelle und Graph von $y = m \cdot x$}
\abhak{1}{Ich kann zu einer Gleichung $y = m \cdot x$ eine Wertetabelle ausfüllen.}
\abhak{2}{Ich kann die Gerade zu $y = m \cdot x$ zeichnen.}
\vgruppe{Werte ablesen und berechnen}
\abhak{3}{Ich kann Werte am Graphen einer proportionalen Funktion ablesen und berechnen.}
\vgruppe{Proportionale Funktion erkennen}
\abhak{4}{Ich kann entscheiden, ob eine Tabelle oder ein Text zu einer proportionalen Funktion gehört.}
\abhak{5}{Ich kann entscheiden, ob eine Tabelle oder ein Text zu einer proportionalen Funktion gehört – weiter.}
\vgruppe{Werte mit dem Dreisatz kontrollieren}
\abhak{6}{Ich kann einen Funktionswert mit dem Dreisatz kontrollieren.}
\vgruppe{Prüfen, begründen, anwenden}
\abhak{7}{Ich finde den Fehler beim Bestimmen des Faktors.}
\abhak{8}{Ich kann begründen, ob eine Gerade zu einer proportionalen Funktion gehört.}
\abhak{9}{Ich kann eine Sachaufgabe zur proportionalen Funktion lösen.}
\abhakgruppe{Einheit 2 · Lineare Funktion $f(x) = m \cdot x + n$}
\abhak{10}{Ich erkenne $m$ und $n$ in der Gleichung.}
\abhak{11}{Ich erkenne ohne Zeichnen, ob eine Gerade steigt oder fällt.}
\abhak{12}{Ich kann an einer Geraden ein Steigungsdreieck einzeichnen.}
\vgruppe{Gerade aus der Gleichung zeichnen}
\abhak{13}{Ich kann zu $f(x) = m \cdot x + n$ eine Wertetabelle anlegen und die Gerade zeichnen.}
\abhak{14}{Ich kann eine Gerade mit $n$ und Steigungsdreieck zeichnen.}
\abhak{15}{Ich kann eine Gerade mit $n$ und Steigungsdreieck zeichnen – weiter.}
\vgruppe{Geraden ablesen und zuordnen}
\abhak{16}{Ich kann $n$ und $m$ an einer Geraden ablesen.}
\abhak{17}{Ich kann einer Gleichung die passende Gerade zuordnen.}
\vgruppe{Parameter deuten}
\abhak{18}{Ich kann aus $m$ und $n$ ablesen, wie eine Gerade verläuft.}
\vgruppe{Prüfen, begründen, anwenden}
\abhak{19}{Ich finde den Fehler beim Ablesen der Steigung.}
\abhak{20}{Ich kann begründen, welche Gerade steiler ist.}
\abhak{21}{Ich kann $m$ und $n$ in einer Sachsituation deuten.}
\abhakgruppe{Einheit 3 · Punkte und Werte}
\abhak{22}{Ich kann einen $x$-Wert in die Gleichung einsetzen.}
\vgruppe{Funktionswerte berechnen}
\abhak{23}{Ich kann einen Funktionswert berechnen.}
\vgruppe{Den $x$-Wert zu einem Funktionswert bestimmen}
\abhak{24}{Ich kann den $x$-Wert zu einem Funktionswert berechnen.}
\vgruppe{Punktprobe}
\abhak{25}{Ich kann prüfen, ob ein Punkt auf einer Geraden liegt.}
\vgruppe{Nullstellen und Achsenschnittpunkte}
\abhak{26}{Ich kann die Nullstelle einer Geraden berechnen.}
\abhak{27}{Ich kann die Nullstelle einer Geraden berechnen – weiter.}
\abhak{28}{Ich kann die Schnittpunkte einer Geraden mit den Achsen angeben.}
\abhak{29}{Ich kann Nullstellen und Schnittpunkte am Graphen ablesen.}
\vgruppe{Prüfen, begründen, anwenden}
\abhak{30}{Ich kann Aussagen über eine Gerade prüfen.}
\abhak{31}{Ich finde den Fehler bei Nullstelle und Punktprobe.}
\abhak{32}{Ich kann begründen, ob und wo eine Gerade eine Nullstelle hat.}
\abhak{33}{Ich kann mit Funktionswert und Nullstelle eine Sachaufgabe lösen.}
\abhakgruppe{Einheit 4 · Gleichung bestimmen}
\vgruppe{Gleichung aus dem Graphen}
\abhak{34}{Ich kann die Gleichung einer Geraden aus dem Graphen bestimmen.}
\vgruppe{Gleichung aus Steigung und Punkt}
\abhak{35}{Ich kann die Gleichung aus der Steigung und einem Punkt bestimmen.}
\vgruppe{Gleichung aus zwei Punkten}
\abhak{36}{Ich kann die Gleichung aus zwei Punkten bestimmen.}
\abhak{37}{Ich kann die Gleichung aus zwei Punkten bestimmen – weiter.}
\abhak{38}{Ich kann eine Gerade durch zwei Punkte zeichnen.}
\vgruppe{Schnittpunkt zweier Geraden}
\abhak{39}{Ich kann den Schnittpunkt zweier Geraden berechnen.}
\abhak{40}{Ich kann den Schnittpunkt zweier Geraden berechnen – weiter.}
\vgruppe{Prüfen, begründen, anwenden}
\abhak{41}{Ich finde den Fehler beim Bestimmen einer Geradengleichung.}
\abhak{42}{Ich kann begründen, wie zwei Geraden zueinander liegen.}
\abhak{43}{Ich kann eine Geradengleichung aus einer Sachsituation bestimmen.}
\abhak{44}{Ich kann mit einem Schnittpunkt eine Sachaufgabe lösen.}
\abhakgruppe{Einheit 5 · Anwendungen}
\abhak{45}{Ich erkenne im Text, was am Anfang da ist und was je Einheit dazukommt.}
\vgruppe{Gleichung zu einem Sachverhalt}
\abhak{46}{Ich kann einem Tarif die passende Gleichung zuordnen.}
\abhak{47}{Ich kann zu einem Sachverhalt die Gleichung aufstellen und einen Wert berechnen.}
\abhak{48}{Ich kann zu einer Gleichung eine passende Situation beschreiben.}
\vgruppe{Tarife vergleichen}
\abhak{49}{Ich kann einem Tarif den passenden Graphen zuordnen.}
\abhak{50}{Ich kann zwei Tarife vergleichen und mich entscheiden.}
\abhak{51}{Ich kann Tarife mit Freiminuten berechnen und vergleichen.}
\vgruppe{Prüfen, begründen, anwenden}
\abhak{52}{Ich finde den Fehler beim Aufstellen einer Tarifgleichung.}
\abhak{53}{Ich kann begründen, ob ein Tarif immer günstiger ist.}
\end{abhakseite}
